% Generated draft: requires mathematical and source review.
\documentclass[11pt]{article}
\usepackage[margin=25mm]{geometry}
\usepackage{fontspec}
\setmainfont{Latin Modern Roman}
\usepackage{amsmath,amssymb,mathtools,unicode-math}
\setmathfont{Latin Modern Math}
\usepackage{enumitem,needspace,xurl,hyperref}
\setlength{\parindent}{0pt}
\setlength{\parskip}{5pt}
\setlength{\emergencystretch}{3em}
\newcommand{\N}{\mathbb{N}}
\newcommand{\Z}{\mathbb{Z}}
\newcommand{\Q}{\mathbb{Q}}
\newcommand{\R}{\mathbb{R}}
\newcommand{\C}{\mathbb{C}}
\newcommand{\F}{\mathbb{F}}
\newcommand{\A}{\mathbb{A}}
\newcommand{\PP}{\mathbb{P}}
\newcommand{\E}{\mathbb{E}}
\newcommand{\eps}{\varepsilon}
\newcommand{\dd}{\,\mathrm{d}}
\begin{document}
\section*{UnsolvedMath ID: 20000309}
\textbf{Source problem code:} catalogue entry
\subsection{Definitions and mathematical statement}
For a normal variety $X$ and an effective boundary $D$ with $K_X+D$ $\R$-Cartier, write
\[
 A_E(X,D)=1+\operatorname{ord}_E(K_Y-\mu^*(K_X+D))
\]
for the log discrepancy of a prime divisor $E$ on a birational model $\mu:Y\to X$. At a specified center $x$, the minimal log discrepancy is the infimum of $A_F(X,D)$ over prime divisors $F$ centered at $x$. A divisor computes it when this infimum equals its discrepancy. Ishii's AIM questions ask which exceptional divisors over mildly singular $X$ compute a nonnegative mld for some pair $(X,D)$, and whether every exceptional component of a surface's minimal resolution can do so. The singularity class, center convention, field and permitted boundaries are not fully specified.

The report treats a singular klt surface germ $(X,x)$. Klt means positive log discrepancies for $(X,0)$. Let $\mu:Y\to X$ be its minimal resolution with exceptional components $E_1,\ldots,E_N$, and put
\[
 P_{ij}=-(E_i\cdot E_j),\qquad v_i=\text{valency of }E_i\text{ in the exceptional dual graph}.
\]
Its claimed finite criterion for a prescribed $E_k$ is existence of a real vector $A=(A_i)$ such that
\[
 A_i\geq0\ \text{for every }i,\qquad
 PA\leq2\mathbf1-v\ \text{componentwise},\qquad
 A_k=\min_iA_i.
\]
The report asserts equivalence between feasibility and existence of an effective $\R$-divisor boundary $B$ for which $E_k$ computes $\operatorname{mld}_x(X,B)\geq0$. The intended vector is $A_i=A_{E_i}(X,B)$. Conversely it claims that every feasible vector is realized by distributing $d=2\mathbf1-v-PA$ among general curvettes, smooth transverse curve germs meeting the exceptional components away from their intersections, with suitable boundary coefficients. This includes a claim about all divisors over $x$, rather than merely minimizing among the listed components; the curvette construction is needed to control further blowups.

As a negative surface example, the report uses the tetrahedral quotient singularity labeled $T_7$, with a central fork curve of self-intersection $-3$ and three arms of lengths $1,2,2$, all five arm curves having self-intersection $-2$. It claims the inequalities force uniformly positive discrepancy gaps from the fork to every arm component. Consequently none of those five arm components can compute a nonnegative mld for any effective boundary in the stated class. The actual gap constants and proof are not supplied.

The field and characteristic assumptions for the klt graph, minimal-resolution and general-curvette constructions are absent; the usual characteristic-zero quotient-surface interpretation requires confirmation. The criterion is a supplied partial surface result, not a classification of all exceptional divisors over arbitrary mildly singular varieties.
\subsection{Short English statement}
Determine which exceptional divisors can compute a nonnegative minimal log discrepancy after adding an effective boundary. The report turns the singular-klt-surface case into finite inequalities on the minimal-resolution graph and claims a curvette realization of every feasible vector. Its $T_7$ example reportedly excludes all arm components, giving a negative answer to unrestricted surface-component optimism.
\subsection{Sources}
\begin{itemize}
\item Local registry: \url{data/unsolvedmath/problems.json}, permanent ID $20000309$, source code \texttt{catalogue entry}; both original questions and supplied research summary.
\item AIM workshop ``Numerical invariants of singularities and higher-dimensional algebraic varieties'', Problem 6, attributed to Ishii: \url{https://aimath.org/WWN/singularvariety/problems_final.pdf}. Canonical location \texttt{aim-algebraic-geometry-notes.json}, \texttt{notes[308]}; original remarks and literature field empty. No PDF was opened here.
\item Supplied machine-generated research attempt 1: partial result; problem partially solved at the run and in the registry; rigor metadata ``proved'', novelty confidence low. No authoritative reference or detailed proof for the inverse criterion, curvette construction or $T_7$ infeasibility is supplied. This non-peer-reviewed aid requires independent verification before citation.
\end{itemize}
\end{document}
